\documentclass[10pt,a4paper]{article}
\usepackage[utf8]{inputenc}
\usepackage[T1]{fontenc}
\usepackage[french]{babel}
\usepackage{geometry}
\usepackage{enumitem}
\usepackage{hyperref}
\usepackage{xcolor}
\usepackage{graphicx}
\usepackage{titlesec}
\usepackage{tikz}

\geometry{top=1.0cm,bottom=0.9cm,left=0.95cm,right=0.95cm}
\setlength{\parindent}{0pt}
\renewcommand{\familydefault}{\sfdefault}
\pagestyle{empty}

\definecolor{maincolor}{RGB}{0,102,153}

\hypersetup{colorlinks=true,urlcolor=maincolor}
\setlist[itemize]{leftmargin=1.15em, itemsep=2.5pt, topsep=2pt, parsep=0pt, partopsep=0pt}

\titleformat{\section}
  {\normalsize\bfseries\color{maincolor}}
  {}
  {0em}
  {}[\vspace{0.25em}\titlerule]

\titlespacing*{\section}{0pt}{0.8em}{0.55em}

\begin{document}

% ===== HEADER =====
\begin{minipage}[t]{0.69\textwidth}
{\LARGE \textbf{LICHIR ISMAIL}}\\[4pt]
\textcolor{maincolor}{\textbf{Alternant Bac +5 | Support technique SAV, diagnostic et réparation}}\\[7pt]

7 impasse Van Gogh, 83100 Toulon\\
Permis B\\[3pt]

\href{mailto:ismaillichir501@gmail.com}{ismaillichir501@gmail.com} \;|\;
\href{tel:+33617997760}{+33 6 17 99 77 60}\\
LinkedIn : \href{https://www.linkedin.com/in/ismail-lichir-3858a4222}{linkedin.com/in/ismail-lichir-3858a4222}\\
GitHub : \href{https://github.com/ismail3080}{github.com/ismail3080}
\end{minipage}
\hfill
\begin{minipage}[t]{0.24\textwidth}
\centering
\vspace{-0.05cm}
\begin{tikzpicture}
\clip (0,0) circle (1.45cm);
\node at (0,-0.03) {\includegraphics[width=3.2cm,height=3.2cm,keepaspectratio=false]{profilpdp.png}};
\end{tikzpicture}
\end{minipage}

\vspace{0.35cm}
{\color{maincolor}\rule{\textwidth}{0.8pt}}
\vspace{0.35cm}

% ===== MAIN CONTENT =====
\begin{minipage}[t]{0.63\textwidth}

\section*{Profil}
Issu d'une formation en systèmes embarqués et électronique, je poursuis actuellement un Master en IoT. Je développe une approche technique orientée intégration matérielle et logicielle, diagnostic de pannes, réparation et maintenance de systèmes connectés. Je m'intéresse particulièrement à la robotique, à l’automatisation industrielle et aux objets connectés.

\section*{Expérience Professionnelle \& Stages}

\textbf{Synapse Audiovisuel – Alternant SAV / Support technique}\\
\textit{Mars 2025 – Présent}
\begin{itemize}
\item Support technique, suivi SAV, diagnostic et réparation de matériels audiovisuels professionnels.
\item Installation, configuration et dépannage de systèmes son/lumière, DSP et processeurs audio.
\item Programmation et configuration Crestron, intégration audiovisuelle et domotique.
\item Diagnostic matériel, électronique et logiciel ; suivi des interventions et accompagnement client.
\end{itemize}

\textbf{YAZAKI MEKNES (YMM) – Stage PFE}\\
\textit{Avril – Juin 2024}
\begin{itemize}
\item Développement d'une application de gestion et de suivi des stocks en environnement industriel.
\item Optimisation d’un processus de tri des PVC et analyse des données de production.
\item Collaboration avec les équipes production et maintenance sur des problématiques d’automatisation industrielle.
\end{itemize}

\textbf{TINKIET France (Hybride) – Stage PFE}\\
\textit{Mars – Juin 2023}
\begin{itemize}
\item Conception de circuits imprimés à base d’ESP32.
\item Programmation de systèmes embarqués.
\item Étude et implémentation de solutions Edge Computing.
\end{itemize}

\textbf{LCM – Aïcha – Stage d’observation}\\
\textit{Juillet 2022}
\begin{itemize}
\item Initiation à la maintenance industrielle et au diagnostic d’équipements.
\end{itemize}

\section*{Formation}

\textbf{Master – Expertise des Systèmes d’Information (IoT)}\\
Epitech Marseille \textit{2025–2027}

\textbf{Licence Sciences pour l'ingénieur}\\
Université de Toulon – La Garde \textit{2024–2025}

\textbf{Licence Sciences et Techniques EEA}\\
Faculté des Sciences et Techniques \textit{2023–2024}

\textbf{DUT Systèmes Embarqués}\\
École Supérieure de Technologie de Fès \textit{2021–2023}

\textbf{Baccalauréat Sciences et Techniques}\\
Lycée Moulay Ismail, Meknès \textit{2018–2021}

\end{minipage}
\hfill
\begin{minipage}[t]{0.33\textwidth}

\section*{Compétences Techniques}

\textbf{Automatisme industriel}\\
Siemens, Schneider, TIA Portal, Grafcet, LOGO!, Twido\\
Automates, capteurs/actionneurs, diagnostic

\textbf{Robotique \& systèmes embarqués}\\
ESP32, STM32, Arduino, Raspberry Pi, FPGA\\
C embarqué, IoT, Edge Computing

\textbf{Audiovisuel \& SAV}\\
Diagnostic, réparation, maintenance\\
Son/lumière, DSP, processeurs audio\\
Support technique, suivi client

\textbf{Crestron \& domotique}\\
Programmation/configuration Crestron\\
Contrôle audiovisuel, automatisation, domotique

\textbf{Électronique \& conception}\\
Altium Designer, Proteus, PSIM, MATLAB/Simulink

\textbf{Communication}\\
CAN, LIN, MQTT, LoRa

\textbf{Programmation}\\
C/C++, Python, Java, JavaScript, React, VHDL

\textbf{Systèmes}\\
Linux (Ubuntu, Kali), RTOS, VMware

\section*{Projets}

\textbf{Tri automatique des tomates}\\
Raspberry Pi, TensorFlow, Edge Impulse

\textbf{TimeManager}\\
Node.js, React, TypeScript, SQL

\textbf{Système photovoltaïque intelligent}\\
MPPT, convertisseur Buck, PSIM

\section*{Certifications}

Systèmes embarqués critiques – INSA\\
Digital Electronics – PROTEUS\\
Edge Computing 3.0\\
OOP Python

\section*{Langues}

Arabe (Maternelle)\\
Français (B2/C1)\\
Anglais (Professionnel)

\end{minipage}

\end{document}
